\begin{statementbox}
    \textbf{\textcolor{CStatement}{条件与图形}}

    椭圆 $E:\dfrac{x^2}{a^2}+\dfrac{y^2}{b^2}=1$，其中 $a>0,b>0$ 且 $a\ne b$；
    点 $P(x_0,y_0)$ 在椭圆外，即
    \[
        S=\frac{x_0^2}{a^2}+\frac{y_0^2}{b^2}>1.
    \]
    从 $P$ 向椭圆作两条切线，切点为不同的 $A,B$。连接 $A,B$，线段 $AB$ 叫\textbf{切点弦}。

    \medskip
    \textbf{\textcolor{CStatement}{结论一：切点弦方程}}
    \[
        \boxed{\frac{x_0x}{a^2}+\frac{y_0y}{b^2}=1.}
    \]
\begin{center}
\begin{tikzpicture}[x=1.28cm,y=1.28cm,font=\small,
  every node/.style={inner sep=3pt},line cap=round]

\coordinate (O) at (0,0); \coordinate (A) at (0,2);
\coordinate (B) at (2.88,-.56); \coordinate (P) at (4,2);
\draw[->,gray!65] (-3.4,0)--(4.6,0) node[right] {$x$};
\draw[->,gray!65] (0,-2.4)--(0,2.7) node[above] {$y$};
\draw[CStatement,very thick] (O) ellipse [x radius=3,y radius=2];
\draw[orange!85!black,thick] (-.65,2)--(4.35,2);
\draw[orange!85!black,thick] (2.55,-1.3143)--(4,2);
\draw[CTrap,very thick] (A)--(B);
\fill (O) circle (1.4pt) node[below left] {$O$};
\fill (A) circle (1.7pt) node[above left] {$A$};
\fill (B) circle (1.7pt) node[below right] {$B$};
\fill (P) circle (1.7pt) node[above right] {$P(x_0,y_0)$};

\node[CTrap,fill=BGStatement] at (.8,.7) {$AB$};

\end{tikzpicture}

{\small 蓝色椭圆；橙色切线 $PA,PB$；红色切点弦 $AB$。}
\end{center}

    \textbf{识别与速写：} 椭圆方程的两个平方项 $x^2,y^2$，分别换成 $x_0x,y_0y$；分母和右端的 $1$ 不变。

    \medskip
    \textbf{\textcolor{CStatement}{结论二：点的位置决定直线的意义}}
    \begin{description}[style=nextline, leftmargin=2.8em, labelsep=0.5em, itemsep=0.35em]
        \item[\textbf{当 $S>1$：}] 该直线经过两个不同的切点，是切点弦所在直线。
        \item[\textbf{当 $S=1$：}] $P$ 在椭圆上，同一方程表示椭圆在 $P$ 处的切线；这时没有两个不同切点。
        \item[\textbf{当 $S<1$：}] $P$ 在椭圆内，不能从 $P$ 作实切线；同一代数方程不是切点弦。特别地，若 $P=O$，方程成为 $0=1$，连直线也不表示。
    \end{description}

    \medskip
    \textbf{\textcolor{CStatement}{几何补充：切点弦中点}}

    在 $S>1$ 时，若 $M$ 是 $AB$ 的中点，则
    \[
        \boxed{M\left(\frac{x_0}{S},\frac{y_0}{S}\right)}.
    \]
    因为 $0<\dfrac1S<1$，所以 $M$ 在从椭圆中心 $O$ 到 $P$ 的线段上；此结论可由把椭圆缩放成圆得到。
\end{statementbox}
